% TOP_PROBLEM: {"id": 20003353, "title": "Strict special-point sequences, Galois component weights, and the correct Haar target", "category_id": 20, "category": {"id": 20, "name": "miscellaneous", "display_name": "Miscellaneous", "description": "Problems whose source classification does not fit the main mathematical categories.", "slug": "miscellaneous", "order_index": 20, "created_at": "2026-07-31T15:26:25.671Z"}, "status": "partially_solved"}
% FIRST_POSTED: 2026-10-01
% ATTEMPT_STATUS: unresolved
% !TeX program = lualatex
% Definition and sources only; this file is not an LLM solution attempt.
\documentclass[11pt]{article}
\usepackage[margin=25mm]{geometry}
\usepackage{fontspec}
\setmainfont{Latin Modern Roman}
\usepackage{amsmath,amssymb,mathtools}
\usepackage{unicode-math}
\setmathfont{Latin Modern Math}
\usepackage{xurl,hyperref}
\hypersetup{colorlinks=true,urlcolor=blue}
\setcounter{secnumdepth}{0}
\setlength{\parindent}{0pt}
\setlength{\parskip}{5pt}
\setlength{\emergencystretch}{3em}
\begin{document}
\section*{\texorpdfstring{Strict special-point sequences, Galois component weights, and the correct Haar target}{Strict special-point sequences, Galois component weights, and the correct Haar target}}\label{20003353}
\subsection{Definitions and mathematical statement}
A special point on a Shimura variety is a zero-dimensional special subvariety, equivalently a point represented by a Shimura datum with toric Mumford--Tate group. In the canonical-model and connected-component convention of the source, let $\mathcal O_j$ be the finite Galois orbit of $x_j$ and set
$\mu_j=|\mathcal O_j|^{-1}\sum_{x\in\mathcal O_j}\delta_x$. Assume that each fixed proper special subvariety contains only finitely many $x_j$. The assertion is weak convergence
\[
 \mu_j\Longrightarrow\mu_{\mathrm{Haar}},
\]
where the ambient invariant measure is normalized to total mass one. The condition is not that a single finite exceptional set of indices excludes all special subvarieties simultaneously. The field of Galois action and component measure must follow the cited setup, especially for a disconnected ambient Shimura variety.
\par\smallskip\textit{Formulation and concept references:} \href{https://aimath.org/WWN/measrigid/measrigid.pdf}{[S1]}.

\subsection{Short English statement}
Do Galois orbits of special points become uniformly distributed when the sequence eventually avoids each fixed proper special subvariety?
\subsection{Research attempt: an analytic reduction with a separate escape-of-mass test}
\textbf{Status.} We derive sufficient, quantitative conditions for the desired convergence, and show why averaging over complete Galois orbits cannot be discarded. The arithmetic estimates required by these conditions are not established. The formulation of orbit and component measure is kept as in the primary workshop source [S1].

\subsubsection{1. A compact spectral criterion}
First suppose the relevant connected quotient is a compact smooth Riemannian manifold $M$ of dimension $d$, with its invariant probability measure $\nu$. This is a stated subcase; orbifold and noncompact issues are not hidden in the notation. Let $\phi_0=1,\phi_1,\ldots$ be an orthonormal Laplace eigenbasis, with eigenvalues $0=\lambda_0\le\lambda_1\le\cdots$. Set
\[a_{j,r}=\int_M\phi_r\,d\mu_j\quad(r\ge1),\qquad
D_s(\mu_j)^2=\sum_{r\ge1}(1+\lambda_r)^{-s}|a_{j,r}|^2,\quad s>d/2.\]
The standard compact-manifold Sobolev embedding gives bounded evaluation on $H^s$, so each finite atomic probability measure belongs to $H^{-s}$ and the displayed quantity is finite. For a smooth real function $f=\sum f_r\phi_r$, weighted Cauchy--Schwarz gives
\[\left|\int f\,d\mu_j-\int f\,d\nu\right|
\le D_s(\mu_j)\left(\sum_{r\ge1}(1+\lambda_r)^s|f_r|^2\right)^{1/2}.\tag{1}\]
Thus $D_s(\mu_j)\to0$ suffices. For a Galois orbit, the coefficient is explicitly
\[a_{j,r}=|\mathcal O_j|^{-1}\sum_{x\in\mathcal O_j}\phi_r(x).\]
The missing input is cancellation in these arithmetic sums, not the norm inequality.

A weaker sufficient condition is $a_{j,r}\to0$ for every fixed $r$. Indeed, finite spectral sums approximate smooth functions uniformly; this follows, for example, by convergence in $H^s$ and Sobolev embedding. Smooth functions uniformly approximate continuous functions on compact $M$. Since both measures have mass one, the two approximation errors are at most twice the uniform error. This proves convergence for every continuous test function. It does not claim that fixed-mode convergence by itself supplies a uniform rate in (1).

\subsubsection{2. The noncompact extension requires control of tails}
For a locally compact quotient, suppose first that convergence is proved against all continuous compactly supported test functions and the proposed limit $\nu$ has total mass one. Choose a compactly supported cutoff $0\le\chi\le1$ with $\int\chi\,d\nu>1-\varepsilon$. Vague convergence implies $\int\chi\,d\mu_j>1-2\varepsilon$ eventually. For any bounded continuous $f$,
\[|\mu_j(f)-\nu(f)|\le |\mu_j(f\chi)-\nu(f\chi)|+
\|f\|_\infty\{\mu_j(1-\chi)+\nu(1-\chi)\}.\]
The right side has limsup at most $3\varepsilon\|f\|_\infty$. Letting $\varepsilon\downarrow0$ proves weak convergence. If the vague limit has mass smaller than one, this argument fails exactly at the cutoff step. Therefore identifying some local limiting density without proving its total mass is insufficient.

\subsubsection{3. A concrete test of the orbit hypothesis}
On the modular curve, the points represented by $\tau_j=ij$ for integers $j\ge2$ are special: $\tau_j$ is imaginary quadratic, satisfying $\tau_j^2+j^2=0$. They lie in the standard cusp region at heights tending to infinity and are pairwise distinct in the modular quotient. Every proper special subvariety of this curve is a point, so this sequence eventually avoids each one. Nevertheless the single-point measures $\delta_{\tau_j}$ escape every compact set and do not converge to normalized invariant area.

This is not a counterexample to the target. A complete Galois orbit generally contains many points besides its chosen representative, and the measures in the target average over those points. The example rigorously rules out a proof that uses only genericity of individual representatives and ignores the arithmetic orbit structure.

\subsubsection{4. Adversarial audit and precise missing input}
All spectral arguments above use the compact smooth subcase explicitly. On a noncompact quotient a discrete eigenbasis need not exhaust the spectrum, so transplanting the compact sum without continuous-spectrum terms would be invalid. The cutoff proof avoids that error but still requires arithmetic convergence on compactly supported functions and a probability limit. Neither the genericity condition alone nor growth of orbit cardinality establishes these cancellations or prevents cusp concentration. Those are the unproved steps toward the full Shimura-variety assertion.

\subsection{Sources}
\begin{itemize}
\item[S1]AIM Workshop Participants, Open Problems from the Workshop Emerging Applications of Measure Rigidity, Conjecture 52, 2004.
\url{https://aimath.org/WWN/measrigid/measrigid.pdf}.
\textit{Formulation source.}
\end{itemize}
\end{document}
